\section{Introduction}

The deployment of autonomous AI agents in Security Operations Centres
(SOCs) is accelerating rapidly. Frameworks such as LangGraph~\cite{langgraph2024},
AutoGen~\cite{wu2023autogen}, and CrewAI~\cite{crewai2024} enable LLM-based
agents to perform threat detection, incident triage, and response
recommendation with minimal human supervision. Yet a fundamental gap persists
in our understanding of these systems: \textit{when do they fail, why do they
fail, and how severely?}

The security community has invested heavily in measuring agent
capabilities (detection accuracy, response speed, ATT\&CK coverage) but
much less in the conditions under which agents produce incorrect,
incomplete, or misleading outputs. This asymmetry carries
serious consequences: a false negative allows an intrusion to persist, a
hallucinated threat attribution poisons shared intelligence feeds, and an
agent paralysed by missing tool access delays incident response beyond
regulatory thresholds.

We identify the core reliability challenge as the
\textit{Skill--Knowledge--Capability (SKC) Mismatch}: the gap between what
a security task demands and what an agent can deliver across several
dimensions. To the best of our knowledge, no existing formal framework ties this
capability mismatch to failure risk in a quantifiable, pre-execution way.

\textbf{Gap in existing work.}
Classical dependability theory~\cite{avizienis2004} provides rigorous fault
taxonomies for conventional software but predates LLM agents and has no
concept of epistemic uncertainty or knowledge currency. Hallucination
research~\cite{huang2023hallucination} analyses LLM reasoning failures but
ignores operational context. AgentBench~\cite{liu2023agentbench} benchmarks
agent task performance without modelling capability constraints or failure
conditions.

\textbf{Contribution.}
We present \ARES{} (Agentic Reliability Evaluation for Security systems),
a compact formal framework for estimating reliability risk in agentic AI
before and during cybersecurity task execution. The conference contribution
is narrow:
\begin{enumerate}[noitemsep, leftmargin=*]
  \item A \textbf{five-dimensional SKC mismatch model} that quantifies the
    gap between agent capability and task demand, producing a learnable
    Agent Reliability Score (ARS) usable as a pre-execution screening
    signal.
  \item An \textbf{Epistemic Gap metric} that formalises the hallucination
    precondition as confident ignorance: the interaction between knowledge
    deficit and miscalibrated confidence.
  \item An \textbf{analytical validation} showing that knowledge
    mismatch dominantly predicts hallucination, the ARS threshold
    generalises across threat classes, and the Epistemic Gap provides an
    interpretable decomposition of false-negative risk into knowledge
    and calibration components.
\end{enumerate}
\ARES{} contributes to intelligent systems research by providing a
reliability-oriented representation and decision framework for autonomous
agents operating under incomplete knowledge and capability constraints.
A cascade reliability bound for tightly-coupled multi-agent pipelines and
a temporal-drift model for knowledge ageing are reported as supporting
results.

The remainder of this paper is organised as follows. Section~\ref{sec:related}
surveys related work. Section~\ref{sec:framework} presents the formal SKC
reliability framework. Section~\ref{sec:epistemic} introduces the Epistemic
Gap. Section~\ref{sec:experimental} provides experimental validation.
Section~\ref{sec:discussion} discusses implications, and
Section~\ref{sec:conclusion} concludes.
